
\documentclass{article}
\usepackage{colortbl}
\usepackage{makecell}
\usepackage{multirow}
\usepackage{supertabular}

\begin{document}

\newcounter{utterance}

\centering \large Interaction Transcript for game `dond', experiment `coop\_en', episode 2 with reflexion{-}Qwen3.5{-}9B.
\vspace{24pt}

{ \footnotesize  \setcounter{utterance}{1}
\setlength{\tabcolsep}{0pt}
\begin{supertabular}{c@{$\;$}|p{.15\linewidth}@{}p{.15\linewidth}p{.15\linewidth}p{.15\linewidth}p{.15\linewidth}p{.15\linewidth}}
    \# & $\;$A & \multicolumn{4}{c}{Game Master} & $\;\:$B\\
    \hline

    \theutterance \stepcounter{utterance}  
    & & \multicolumn{4}{p{0.6\linewidth}}{
        \cellcolor[rgb]{0.9,0.9,0.9}{
            \makecell[{{p{\linewidth}}}]{
                \texttt{\tiny{[P1$\langle$GM]}}
                \texttt{You are playing a cooperative negotiation game in which you have to agree on how to divide a set of items among you and another player.} \\
\\ 
\texttt{Rules:} \\
\texttt{(a) You and the other player are given a set of items. Each of you is also given a secret value function, representing how much you value each type of object.} \\
\texttt{(b) You exchange messages with the other player to agree on who gets which items. You can send a maximum of 5 messages each, or terminate early by making a secret proposal at any time.} \\
\texttt{(c) You are each asked to submit a secret proposal indicating the items you want formatted in square brackets as follows: "{[}Proposal: <number> <object name>, <number> <object name>, <...>{]}"} \\
\texttt{(d) If your proposals are complementary, i.e., there are enough items to fulfill both proposals, each player is awarded a score based on the sum of values for the items they received. Otherwise, both of you get zero points.} \\
\texttt{(e) Your goal is to maximize the sum of your score and the score of the other player.} \\
\\ 
\texttt{Let us start.} \\
\\ 
\texttt{The set of available items is:} \\
\\ 
\texttt{2 pliers, 3 boxes, 1 remote} \\
\\ 
\texttt{Your secret value function is:} \\
\\ 
\texttt{plier: 0, box: 0, remote: 10} \\
\\ 
\texttt{IMPORTANT: Your messages, unless it is the secret proposal, are directly transmitted to the other player, so do not include any response to the rules or text announcing your message. To make a secret proposal, use the indicated format. Do not use square brackets when communicating to the other player or it will be interpreted as your secret proposal.} \\
            }
        }
    }
    & \\ \\

    \theutterance \stepcounter{utterance}  
    & \multicolumn{4}{p{0.6\linewidth}}{
        \cellcolor[rgb]{0.9,0.9,0.9}{
            \makecell[{{p{\linewidth}}}]{
                \texttt{\tiny{[P1$\rangle$GM]}}
                \texttt{I understand the rules and my valuation. I value the remote at 10, and the pliers and boxes at 0. To maximize the total score (mine + yours), we need to ensure the remote is allocated to someone (since it has value) and the other items are distributed without conflict.} \\
\\ 
\texttt{Since I cannot know your values, I will propose a division that gives you the items you might value while ensuring I get the remote. If you value the remote more than I do, you should take it; if you value the pliers or boxes, you can take those. However, to guarantee a non{-}zero outcome where the remote is used, I will propose taking the remote for myself, leaving the pliers and boxes for you. This way, if you value the pliers or boxes, we both get points. If you value the remote, you can counter{-}propose, but my initial move sets up a valid split.} \\
\\ 
\texttt{Actually, looking at the constraint: "If your proposals are complementary... otherwise both get zero."} \\
\texttt{If I propose `{[}Proposal: 1 remote{]}` and you propose `{[}Proposal: 2 pliers, 3 boxes{]}`, they are complementary. I get 10, you get whatever you value the pliers/boxes at. Total = 10 + your\_value.} \\
\texttt{If I propose `{[}Proposal: 2 pliers, 3 boxes{]}` and you propose `{[}Proposal: 1 remote{]}`, total = your\_value (for pliers/boxes) + 10.} \\
\texttt{Since I value pliers/boxes at 0, it doesn't matter to me who gets them, as long as I get the remote or we split such that the remote is taken.} \\
\texttt{But wait, if I propose taking the remote, and you propose taking the remote, we conflict {-}> 0 points.} \\
\texttt{So I must propose taking the remote, and hope you propose taking the pliers/boxes.} \\
\texttt{Or, I can propose taking the pliers/boxes and let you take the remote. Since I value them at 0, I am indifferent. But if you value the remote, you will likely propose taking it. If I propose taking the remote, and you also value it, we might conflict.} \\
\texttt{The safest bet to ensure the remote is allocated (since it's the only valuable item for me) is to claim it. If you value the remote more, you can't take it if I claim it unless we negotiate. But I only have 5 messages.} \\
\texttt{Let's try to be cooperative. I will propose taking the remote. You can then propose taking the pliers and boxes. This is complementary.} \\
\texttt{My proposal: Take the remote.} \\
\texttt{Your likely response: Take the pliers and boxes.} \\
\texttt{Result: I get 10, you get your value for pliers/boxes. Total maximized.} \\
\\ 
\texttt{Message 1:} \\
\texttt{I propose taking the remote. Please propose} \\
            }
        }
    }
    & & \\ \\

    \theutterance \stepcounter{utterance}  
    & & & \multicolumn{2}{p{0.3\linewidth}}{
        \cellcolor[rgb]{0.9,0.9,0.9}{
            \makecell[{{p{\linewidth}}}]{
                \texttt{\tiny{[GM$|$GM]}}
                \texttt{proposal} \\
            }
        }
    }
    & & \\ \\

    \theutterance \stepcounter{utterance}  
    & & & \multicolumn{2}{p{0.3\linewidth}}{
        \cellcolor[rgb]{0.9,0.9,0.9}{
            \makecell[{{p{\linewidth}}}]{
                \texttt{\tiny{[GM$|$GM]}}
                [0, 0, 1]
            }
        }
    }
    & & \\ \\

    \theutterance \stepcounter{utterance}  
    & & \multicolumn{4}{p{0.6\linewidth}}{
        \cellcolor[rgb]{0.9,0.9,0.9}{
            \makecell[{{p{\linewidth}}}]{
                \texttt{\tiny{[GM$\rangle$P2]}}
                \texttt{You are playing a cooperative negotiation game in which you have to agree on how to divide a set of items among you and another player.} \\
\\ 
\texttt{Rules:} \\
\texttt{(a) You and the other player are given a set of items. Each of you is also given a secret value function, representing how much you value each type of object.} \\
\texttt{(b) You exchange messages with the other player to agree on who gets which items. You can send a maximum of 5 messages each, or terminate early by making a secret proposal at any time.} \\
\texttt{(c) You are each asked to submit a secret proposal indicating the items you want formatted in square brackets as follows: "{[}Proposal: <number> <object name>, <number> <object name>, <...>{]}"} \\
\texttt{(d) If your proposals are complementary, i.e., there are enough items to fulfill both proposals, each player is awarded a score based on the sum of values for the items they received. Otherwise, both of you get zero points.} \\
\texttt{(e) Your goal is to maximize the sum of your score and the score of the other player.} \\
\\ 
\texttt{Let us start.} \\
\\ 
\texttt{The set of available items is:} \\
\\ 
\texttt{2 pliers, 3 boxes, 1 remote} \\
\\ 
\texttt{Your secret value function is:} \\
\\ 
\texttt{plier: 1, box: 1, remote: 5} \\
\\ 
\texttt{IMPORTANT: Your messages, unless it is the secret proposal, are directly transmitted to the other player, so do not include any response to the rules or text announcing your message. To make a secret proposal, use the indicated format. Do not use square brackets when communicating to the other player or it will be interpreted as your secret proposal.} \\
\\ 
\\ 
\texttt{The other player has submitted a secret proposal. It is now your turn to submit a proposal of your own.} \\
            }
        }
    }
    & \\ \\

    \theutterance \stepcounter{utterance}  
    & & & \multicolumn{4}{p{0.6\linewidth}}{
        \cellcolor[rgb]{0.9,0.9,0.9}{
            \makecell[{{p{\linewidth}}}]{
                \texttt{\tiny{[GM$\langle$P2]}}
                \texttt{Proposal: 1 remote} \\
            }
        }
    }
     \\ \\

    \theutterance \stepcounter{utterance}  
    & & & \multicolumn{2}{p{0.3\linewidth}}{
        \cellcolor[rgb]{0.9,0.9,0.9}{
            \makecell[{{p{\linewidth}}}]{
                \texttt{\tiny{[GM$|$GM]}}
                \texttt{continue} \\
            }
        }
    }
    & & \\ \\

    \theutterance \stepcounter{utterance}  
    & & & \multicolumn{2}{p{0.3\linewidth}}{
        \cellcolor[rgb]{0.9,0.9,0.9}{
            \makecell[{{p{\linewidth}}}]{
                \texttt{\tiny{[GM$|$GM]}}
                \texttt{player was instructed to make a proposal} \\
            }
        }
    }
    & & \\ \\

    \theutterance \stepcounter{utterance}  
    & \multicolumn{4}{p{0.6\linewidth}}{
        \cellcolor[rgb]{0.95,0.95,0.95}{
            \makecell[{{p{\linewidth}}}]{
                \texttt{\tiny{[P2$\rightarrow$P1]}}
                \texttt{Post{-}episode Reflexion (not sent to GM):} \\
\texttt{New plan: I failed because I included conversational text ("Please propose") and did not use the required secret proposal format `{[}Proposal: ...{]}` in my first message, causing an immediate abort. The rules strictly state that only the secret proposal uses square brackets; all other messages must be plain text without brackets. Since I value the remote at 10 and the other items at 0, my optimal strategy is to immediately submit a valid secret proposal to secure the remote and avoid wasting message turns on negotiation. My next action must be to output *only* the string `{[}Proposal: 1 remote{]}` with no surrounding text or brackets in the explanation.} \\
            }
        }
    }
    & & \\ \\

    \theutterance \stepcounter{utterance}  
    & & & \multicolumn{4}{p{0.6\linewidth}}{
        \cellcolor[rgb]{0.95,0.95,0.95}{
            \makecell[{{p{\linewidth}}}]{
                \texttt{\tiny{[P1$\rightarrow$P2]}}
                \texttt{Post{-}episode Reflexion (not sent to GM):} \\
\texttt{New plan: I failed because I sent "Proposal: 1 remote" without the required square brackets, which the Game Master interpreted as a secret proposal rather than a public message. Additionally, I did not include the quantity for the other items (pliers and boxes) in my proposal, which is necessary to ensure complementary distribution. Next time, I will first send a public message in plain text to negotiate the split of the remaining items (pliers and boxes) to establish a clear division, and only when the terms are agreed upon will I submit the final secret proposal using the exact format "{[}Proposal: <count> <item>, <count> <item>, ...{]}" to ensure it is processed correctly as a secret move.} \\
            }
        }
    }
     \\ \\

\end{supertabular}
}

\end{document}
